\documentclass[10pt]{article}
\usepackage[margin=20mm]{geometry}
\usepackage{amsmath,amssymb,mathtools,fontspec,unicode-math,xurl,hyperref}
\IfFontExistsTF{DejaVu Serif}{\setmainfont{DejaVu Serif}}{\setmainfont{Latin Modern Roman}}
\setmathfont{latinmodern-math.otf}
\setlength{\parindent}{0pt}\setlength{\parskip}{6pt}\sloppy
\title{Technical verification of bounded-type Herman-ring estimates}
\author{AI-assisted mathematical review}
\date{}
\begin{document}
\maketitle
Abstract: This report examines the exact author manuscript arXiv:2302.07794v5 of Willie Rush Lim. An almost-tiling selection is impossible as literally stated, the initial-level construction of first-return domains degenerates, and a scale cutoff is off by one. Explicit bounded corrections preserve the quantitative conclusions examined. A prescribed-return-time transition and a geometric width comparison require proof supplements, which are distinguished from their original formulations. No counterexample to the principal theorems was established. The impossible selection is a Major defect in the core amplification proof; the other two findings have bounded local impact. Verified repairs and the remaining dependency limits do not certify the complete paper. Scientific adjudication remains pending.\par
\section*{Identity, source, and scope}
The paper is Willie Rush Lim, A priori bounds and degeneration of Herman rings with bounded type rotation number, arXiv:2302.07794v5, 9 September 2025, 59 PDF pages, [Lim], associated with DOI 10.1007/s00222-025-01369-2 (https://doi.org/10.1007/s00222-025-01369-2). All locations below refer to this exact author’s PDF. Its printed page numbers agree with the PDF page indices. The journal’s typeset version was not obtained; neither textual equivalence nor the persistence of these observations in that version is asserted.\par
The available record comprises a complete reading of the source and bounded checks of its extremal-width, amplification, compactness, realization, and first-return arguments. This report integrates those preserved checks with targeted recovery of the outstanding proofs and page evidence. The coverage ledger distinguishes reading, verified derivation, externally invoked results, and incomplete independent reconstruction. The external results cited below were checked only in the documented relevant ranges; their complete foundational proofs were not reproduced. No author code was run. Missing author code is irrelevant to the mathematical assessments made here.\par
\section*{Definitions and operative hypotheses}
The principal family has fixed local degrees \(d_0,d_\infty\geq2\) at the superattracting fixed points \(0,\infty\), total degree \(d=d_0+d_\infty-1\), and an invariant separating Herman ring of rotation number \(\theta\). Every remaining critical point lies on its two boundaries. The rotation is of bounded type; write \(B=\sup_j a_j<\infty\) for its continued-fraction bound. Theorems for this family are not statements about arbitrary rational Herman rings.\par
Use the source conventions \(q_{-1}=0\), \(q_0=1\), \(q_n=a_nq_{n-1}+q_{n-2}\) and \(l_n=|q_n\theta-p_n|\). Lengths on the distinguished invariant set are measured in the rotation coordinate. Constants described as absolute may depend on \(d_0,d_\infty,B\). A quasicircle, a closed thick ring, and a boundary circle at the Siegel scale are different geometric cases. In particular, the latter scale requires the interval length to be no greater than the ring modulus.\par
The extremal width \(\mathcal W\) is the reciprocal of extremal length. For canonical horizontal types \(A,B\) the weighted quantity is \(Z=\mathcal W(A)+2\mathcal W(B)\), so \(\mathcal W\leq Z\leq2\mathcal W\). For positive numbers write \(a\oplus b=ab/(a+b)\). The distinction between total covering degree and local degree over a marked collar is essential: thresholds can depend on the former, whereas the multiplicative gain used for uniform amplification must use the latter.\par
\section*{An impossible narrow-gap selection}
Problem 1 (Almost-tiling pieces cannot always satisfy the stated gap window).  In the deep-case proof of Theorem 5.1, pp. 23–24, the source chooses \(N=\lfloor\lambda/(3\tau^2)\rfloor\) level-\(n\) almost-tiling pieces with consecutive gaps in \(((\tau-1)l_n,\tau l_n]\), where \(\tau=10\). The same narrow-window selection is used in Lemma 6.2, p. 28. For golden-mean rotation and every even \(n\geq8\), a chain meeting this literal condition contains at most four pieces.\par
\subsection*{Exact certificate and its assumptions}
Put \(\theta=(\sqrt5-1)/2\) and \(\varphi=1/\theta\). For even \(n\), let \(q=q_n\), \(p=q_{n-1}\), \(Q=q+p=q_{n+1}\), and \(r=p/q\). The allowed initial points are \(c+i\theta\), \(0\leq i<Q\), and the pieces have positive angular length \(l_n\). Translation of \(c\) changes none of the following differences. The recurrence gives \[\frac35\leq r\leq\frac58,\qquad
  q_n\theta-p_n=l_n,\qquad
  q_{n-1}\theta-p_{n-1}=-\varphi l_n.\] The interval for \(r\) is invariant under the two-step update \(r\mapsto(1+r)/(2+r)\).\par
A permitted forward difference of initial points lies in \((10l_n,11l_n]\). If \(k\) is its time-index difference, then \(|k|<Q\). Unimodularity of consecutive convergents gives unique integers \(a,b\) with \[(k,m)=a(q_n,p_n)+b(q_{n-1},p_{n-1}),\quad
 E=a-b\varphi\in(10,11],\quad
 k/q=E+b(\varphi+r).\] If \(b\geq-3\), this last expression is greater than \(13/4>1+r\); if \(b\leq-6\), it is less than \(-11/5<-(1+r)\). Thus \(b=-4\) or \(-5\). The respective integer values are \(a=4\) and \(2\). There are exactly two possible steps, \[A=4q-4p,\qquad B_*=2q-5p.\] The admissible starting bands are \[L_+=[0,-3q+5p-1],\qquad L_-=[-2q+5p,Q-1].\] Step \(A\) sends \(L_+\) into \(L_-\). An \(A\) step followed by \(B_*\) lands in \([6q-9p,3q-4p-1]\), strictly between the two starting bands. Step \(B_*\) cannot land in \(L_-\). Consequently every chain has at most three steps, hence four pieces.\par
This conclusion also covers circular shortest-distance readings. A large forward difference would be at least \(1-11l_n\). A five-piece ordered chain can have at most one such gap; that gap and the other three differences, each greater than \(10l_n\), would total more than one. Thus the large-gap alternative does not evade the certificate. Endpoint conventions do not help, since both permitted errors are strictly inside the specified interval.\par
This is an analytic argument at arbitrarily deep even levels, not extrapolation from numerical experiments. For any fixed arbitrarily large \(\lambda\) with \(N\geq5\), \(n\) can be chosen beyond the finite deep-case thresholds. The independently written exact finite checks at \(n=14,20,28\) agree with it. For example, \(q=514229\), \(p=317811\) at \(n=28\), and the maximal chain \(560597\to0\to785672\to225075\) has only four vertices. Full exact arithmetic and the parameter-binding certificate are retained as evidence E1.\par
\subsection*{Counterevidence, correction, and impact}
This certificate refutes the intermediate combinatorial selection. It does not exhibit a rational map satisfying an arbitrarily large-width hypothesis and therefore is not a counterexample to Theorem 5.1.\par
A greedy replacement chooses the first subsequent piece with gap greater than \((\tau-1)l_n\). Adjacent initial-point separations are less than \(2l_n\), so the new upper gap bound is \((\tau+1)l_n\). The selected cluster has length \(L<(\tau+2)Nl_n\). The separation of the \(\tau\)-neighborhoods is unchanged. Moreover \(\tau P\subset\lambda I_j\) follows from \((\tau+1)L\leq\lambda l_n\), because at \(\tau=10\) \[(\tau+1)(\tau+2)N\leq 132\lambda/300<\lambda .\] Thus the same \(N\) and the same quasi-additivity gain \(\sqrt{N/2}\,K\) are available. The level offset remains \(O_B(\log\lambda)\), with no dependence on the modulus. The analogous fixed constants in Lemma 6.2 can be enlarged without changing their independence from \(\lambda\).\par
The supported assessment is Major. Theorem 6.1 and Proposition 4.2 feed the deep-case gain of Theorem 5.1; the impossible selection is an input to its Quasi-Additivity argument. That gain is then used in Theorem 8.1, Corollary 8.2 and the proof of Theorem A (pp. 44–46). This is an invalid derivation in a core proof. The verified wider-window repair retains the same gain and parameter dependence; repair status is recorded separately and no counterexample to Theorem A is asserted. The selection should be corrected at both source locations.\par
\section*{The initial first-return construction}
Here \(f=f_\theta\) is a unicritical boundary realization selected by Theorem C. Its unique free critical point is \(c_0=1\), with local degree \(d\), and \(c_j=f^j(c_0)\) lies on the invariant curve. These additional hypotheses, rather than the general multicritical family alone, justify the single-cycle covering argument below.\par
Problem 2 (The displayed enlarged interval degenerates at initial levels).  Theorem 10.2, p. 51, states a pointed first-return construction for every \(n\). Its proof uses enlarged-interval endpoints \(c_{q_n-q_{n-1}}\) and \(c_{q_{n-1}}\). For golden-mean combinatorics, \(q_0=q_1=1\), \(q_2=2\). At \(n=2\) both endpoints are \(c_1\); at \(n=1\) the proposed target \(c_{q_1}\) is an endpoint. The displayed construction therefore does not provide the asserted nondegenerate pointed collar at these levels.\par
These are exact integer equalities under the source’s own recurrence, verified against pp. 13 and 51. They are not extraction errors, roundoff, or an allegation that the stated first-return map cannot exist. A shortest-arc ambiguity can also occur at \(n=3\). For \(n\geq4\) the intended enlarged arc has length \(2l_{n-1}+l_n<1/2\), so that finite-prefix ambiguity disappears. Evidence E2 gives the full arithmetic and source-page binding.\par
Here is a repair with the required uniform dependence. For \(n\leq3\), put \(Q_0=(B+1)^3\), so \(q_n\leq Q_0\). The bounded-type estimate \[\|k\theta\|_{\mathbb R/\mathbb Z}\geq((B+2)|k|)^{-1}\] uniformly separates the finitely many relevant orbit points. Use the normalized uniform quasiconformal conjugacy supplied by the limiting family, fixing \(0,\infty\) and taking the invariant curve to the unit circle. Choose two concentric disks about the image of \(c_{q_n}\) with small radii \(r,2r\) depending only on \(B\). Their inverse images \(V,\widetilde V\) avoid all other critical values of \(f^{q_n}\) and have a uniform collar.\par
The component \(\widetilde U\) of the preimage containing \(c_0\) has degree \(d\): the only critical-value monodromy in the target disk is the \(d\)-cycle at \(c_0\); connectedness excludes extra fixed sheets in that component. Riemann–Hurwitz gives a disk. In Riemann coordinates the map is a rotated \(z^d\). Pulling back the smaller target disk and applying Koebe distortion yields bounded pointed shape. Uniform quasisymmetry and finite-prefix separation give the required two-sided comparison with \(|c_0-c_{q_n}|\). The full argument, including the normalization that keeps \(0,\infty\) separated, is retained in E2.\par
The assessment is Minor: the defect concerns a finite prefix of the specific construction. The repaired statement retains the uniform dependence, and the subsequential rescaled limit in Corollary 10.3 is unaffected by deleting a finite prefix. No existence or compactness theorem is refuted. For the latter corollary, degree preservation was checked separately by writing each normalized map in Riemann coordinates as \(e^{i\alpha_n}z^d\); bounded pointed domains give nonconstant limiting Riemann maps. Bounded shape alone was not treated as a general no-degree-loss theorem.\par
\section*{A boundary between the two scale regimes}
Problem 3 (The Siegel-scale range starts one level too early).  Equation (8.1), p. 45, defines \(M\) by \(l_{M+1}\leq\mu<l_M\), but the subsequent text treats every \(n\geq M\) as Siegel scale. At \(n=M\) the interval length is \(l_M>\mu\), contrary to the earlier scale hypothesis.\par
This is a concrete mismatch of applicability, rather than a numerical counterexample to the final a priori bound. Set \(S=M+1\). Then the Herman range ends at \(n<S\) and the Siegel range starts at \(n\geq S\). In Lemma 8.3 the undefined cutoff \(N_1\) should likewise be replaced by the actual threshold \(S\). The passage through the transition band is then \(S\leq n_1<S+tm\). Since \(\mu<l_{S-1}\), the same bounded-type estimates give \(\mu/l_{n_1+1}\leq C^{tm+1}\). The fixed transition losses and the order of choices are preserved. Evidence E3 verifies pp. 45–46 and this substitution.\par
The assessment is Minor, supported by the one-level range correction and unchanged quantitative transition argument. Isolated incorrect interval labels elsewhere were not counted as additional substantive findings.\par
\section*{Proof supplements and unresolved limits}
\subsection*{From a stopping time to the specified denominator}
Proposition 7.2, p. 33, requests small horizontal width at the prescribed time \(q_{n+m_2}\). Its printed proof on p. 43 produces a time \(a=tq_{n+m}\), \(1\leq t\leq\bar t\). A small width at that stopping time does not by itself state the prescribed-time bound; a multiple of a convergent denominator need not be a convergent denominator.\par
For the proper disk covers and disjoint full marked continua used in Section 7, the finite-cover middle-cut argument supplies the needed supplement under its actual hypotheses. If \(a+b\leq q_{n+k}\) for fixed \(k\), the marked continua have local degree at least two, and the finite topology is controlled, it gives \[Z^{a+b}\leq Z^a+C(k,\lambda).\] The type-\(A\) once/type-\(B\) twice crossing argument, peripheral rectangles, and bounded canonical-restriction loss must be retained. For sufficiently large fixed \(k_0\), \(b\geq q_{n+k_0}\) ensures the required precritical encounters in the marked intervals. Choose the stopping offset \(m\geq k_0\) at the outset, and run the stopping argument with \(\eta=\delta/8\). Choose \(m_2\) with \(q_{n+m_2}\geq(\bar t+1)q_{n+m}\), using \(q_{j+2}\geq2q_j\). Then \(b=q_{n+m_2}-a\) is large enough, and \[\mathcal W(F^{q_{n+m_2}})\leq Z^a+C\leq2\eta K+C<\delta K\] after the permitted enlargement of \(K\). The order is \(\Delta,\delta\), then \(\eta,M,m,\nu,\bar t,m_2\), then \(C(\lambda,m_2)\), then \(K\). It preserves independence of \(m_1,m_2\) from \(\lambda\).\par
This is a proof-completeness clarification with a bounded supplement, not a verified false proposition. The step feeds Lemma 6.6 and Theorem 6.1 and should be written explicitly. Evidence E3 distinguishes the original stopping proof from the reconstructed bridge. No unrestricted monotonicity of widths under arbitrary branched maps is being assumed.\par
\subsection*{The literal domain inclusion and the inequality actually used}
The cut construction on p. 30 does not by itself certify the literal inclusion (6.3). No rational-map counterexample to that inclusion was obtained. For the required p. 32 width inequality, a direct replacement suffices: the new full, disjoint marked continua are compactly contained in the new disk and contain the old ones. A type-\(A\) leaf has an essential horizontal or vertical first-transition segment. A type-\(B\) leaf that hits the opposite enlarged continuum or leaves the new disk also has such a first-transition segment. Otherwise, nonzero winding in the complement of the opposite marked continuum forces a segment between different lifted copies of its own marked continuum. Overflow, the parallel law, and canonical maximality then give the desired horizontal-plus-vertical width bound with an absolute topological loss. This argument does not assert the disputed literal inclusion. Its full hypotheses and proof are in E3.\par
The geometric canonical-restriction statement in Appendix A.12 was also checked separately from the wording of Kahn’s weighted-arc-diagram corollary [Kahn]. A strip metric bounds the discarded bad-segment width by four per noncrossing arc class, giving a topology-dependent total loss. The stronger geometric statement was not silently attributed to a weaker external formula.\par
\subsection*{A dependency that remains only partially reconstructed}
The additional peripheral \(E\)-absorption and canonical reidentification in Section 7.1 was read in context, but its complete surgery with quantitative buffers was not independently rederived. The established restriction lemma does not, by itself, prove that additional operation. A separate stronger outer-buffer assertion on pp. 19–20, retaining a subfamily of an arbitrary original lamination, also remains uncertified: the checked rectangular-wave replacement supplies the width needed downstream but does not establish that its new vertical leaves belong to the original family. No counterexample to that stronger wording was constructed. These are limits of the audit, not confirmed mathematical defects. The principal conclusions also depend on external quasiconformal, realization, and extremal-length theorems that were invoked with checked hypotheses, rather than re-proved in their entirety.\par
\section*{Positive checks and overall dependence}
The recovered analytic checks support several potentially delicate parts of the argument. In particular, \[\bigl((a+c)\oplus(b+e)\bigr)-(a\oplus b)-(c\oplus e)
 =\frac{(ae-bc)^2}{(a+b)(c+e)(a+b+c+e)}\geq0\] for positive parameters, verifying the aggregation direction in (7.2). The scalar contraction in Lemma 7.14 satisfies \[1-\nu'=\frac{81}{10(20\kappa+29)}>0,\qquad\kappa>1 .\] These are exact identities, not empirical checks.\par
The global amplification choices are noncircular: fix the gain \(\chi\) first, then \(\Xi=2/\chi\) and its spreading constant, then \(\lambda\), and finally the large-width threshold. The fixed-map finite width bound is sufficient for the contradiction proving a uniform threshold; uniformity is not assumed at the start. The thick-ring transition preserves this logic with a finite band of levels.\par
For the degeneration argument, zero–pole separation prevents loss of degree. Uniform attracting neighborhoods exclude the free Julia critical points from colliding with \(0\) or \(\infty\); collisions among free critical points are allowed. For a zero-modulus limit, the postcritical set is the invariant curve together with the two attractors. The rotation-domain boundary theorem [McMullen] then excludes an additional rotation domain without assuming that its rotation is of bounded type. The realization input uses an explicit reviewer-constructed topological model together with the checked obstruction and surgery hypotheses of Zhang, Wang, and Shishikura [Zhang2, Wang, Shishikura]; Wang’s obstruction criterion alone was not quoted as arbitrary-tuple existence. Zhang’s quantitative proof was traced to obtain the required uniform dependence on a bound for the continued-fraction entries, instead of silently changing a literal \(K(d,\theta)\) statement to \(K(d,B)\) [Zhang]. The relevant Log Rule and local-versus-total-degree collar estimates were also checked in their primary formulation [DL].\par
Finally, for the unicritical family with \(m=d_0\), \(n=d_\infty\), \(d=m+n-1\), set \[Q(z)=\sum_{j=0}^{m-1}(-1)^j{d\choose j}z^j,\qquad
 P(z)=\sum_{j=m}^{d}(-1)^j{d\choose j}z^j .\] Since \(P+Q=(1-z)^d\) and \(Q(1)=(-1)^{m-1}{d-1\choose m-1}\ne0\), the formula \(F_c=-cP/Q\), \(c\in\mathbb C^*\), has no cancellation. The identity \[dQ+(1-z)Q'=(-1)^{m-1}\frac{d!}{(m-1)!(n-1)!}z^{m-1}\] gives the exact critical pattern and ramification total \(2d-2\). This does not imply that every \(c\) has a Herman curve. Conjecture D and parameter-plane illustrations were not treated as proved theorems.\par
\section*{Conclusion}
The supported paper-level assessment is Major: one core proof defect and two Minor findings are established, with explicit verified repairs, proof-completeness clarifications and the stated incomplete dependency audit. This judgment concerns the exact inspected author manuscript and the demonstrated effects. It neither requires nor asserts that the whole paper has been refuted.\par
Question | Recorded outcome\par
Narrow-gap selection | Major core proof defect; bounded replacement checked\par
Initial first-return domains | Finite-prefix defect established; uniform repair checked\par
Siegel cutoff | One-level applicability mismatch; transition correction checked\par
Prescribed-time bridge | Proof supplement checked under stated geometric hypotheses\par
Literal inclusion (6.3) | Not established or refuted; needed inequality proved separately\par
Peripheral absorption | Independent reconstruction incomplete; no defect alleged\par
Scientific adjudication | Pending\par
Independent acceptance | Current focused checks recorded in verification\par
Human review / author adoption | No human review supplied / adoption unknown\par
\section*{Supporting checks}
All evidence paths are relative to the assigned staging root. They are retained with actual byte counts and SHA256 bindings in the verification records; the old execution paths are preserved separately from the new paths.\par
\paragraph*{E1.}\mbox{}\par
The source pages used include 23–24 and 28.\par
\paragraph*{E2.}\mbox{}\par
Relevant source pages include 13 and 50–52.\par
\paragraph*{E3.}\mbox{}\par
Relevant source pages include 30–37, 43, and 45–46.\par
\paragraph*{Other checks.}\mbox{}\par
Historical checks are reused only with their actual bound inputs. New displays are distinguished from historical displays; rendering by itself is never counted as inspection.\par
\paragraph*{Supporting checks}\mbox{}\par
The source-manifest SHA256 is `d09595bba91df12b78dcaa016d1dec296532e7e399a778d5e6d6ae3062eb0bdd`. Original and recovery scientific work retains its own conversation and native-helper attribution. The external controller’s UI observation and the runtime instruction self-description are recorded separately in provenance; neither is an independent measurement of a backend model. An earlier operational role-audit helper supplied no mathematical evidence. No institutional endorsement or human review is claimed.\par
The final source, compiler, input/dependency, command, exit-code, output, render, and actual all-page inspection records are attached in the build and provenance files. They are mechanical provenance, not scientific acceptance.\par
9 W. R. Lim, A priori bounds and degeneration of Herman rings with bounded type rotation number, arXiv:2302.07794v5, 9 September 2025. J. Kahn, A priori bounds for some infinitely renormalizable quadratics: I. Bounded primitive combinatorics, arXiv:math/0609045v2 (2006) and v3 (2023). Only the separately identified canonical-arc results were used. D. Dudko and M. Lyubich, Uniform a priori bounds for neutral renormalization, arXiv:2210.09280v1, 2022. G. Zhang, All bounded type Siegel disks of rational maps are quasi-disks, arXiv:0810.2733v4, 2010. C. T. McMullen, Complex Dynamics and Renormalization, author-posted copy; Corollary 3.7 and the specified material in Sections 5.1–5.2. X. Wang, A decomposition theorem for Herman maps, arXiv:1203.5563v1; version watermark and differing body date retained in provenance. G. Zhang, Dynamics of Siegel rational maps with prescribed combinatorics, arXiv:0811.3043v1, 2008. M. Shishikura, On the quasiconformal surgery of rational functions, Ann. Sci. École Norm. Sup. (4) 20 (1987), 1–29; Sections 3 and 9.\par
\end{document}
